% !PRINTED 731
\weberpage{731}

Hieraus kann, wie im \S.\,187, gefolgert werden, dass es Systeme von $\nu'$ unabhängigen Normaleinheiten giebt, deren Regulator $T_0$ \gesperrt{so klein als möglich} wird, und solche Systeme heissen \gesperrt{Fundamentalsysteme von Normaleinheiten}.

Ist $\mathcal{E}_0, \mathcal{E}_1, \ldots, \mathcal{E}_{\nu'-1}$ ein solches Fundamentalsystem, so ergeben sich die Exponenten $\xi_0, \xi_1, \ldots, \xi_{\nu'-1}$ in (10) und (11) als \gesperrt{ganze rationale Zahlen}. Auch dies kann ebenso bewiesen werden, wie im \S.\,187.

Stellen wir unter dieser Voraussetzung über die $\mathcal{E}_i$ die Gleichungen (11) für irgend ein anderes System unabhängiger Einheiten $\mathcal{E}_0^{(1)}, \mathcal{E}_1^{(1)}, \ldots, \mathcal{E}_{\nu'-1}^{(1)}$ auf, so erhalten wir Ausdrücke von der Form:

\[
L_{i,k}^{(1)} = \xi_{i,0} L_{0,k} + \xi_{i,1} L_{1,k} + \cdots + \xi_{i,\nu'-1} L_{\nu'-1,k}.
\]

Ist also $T_1$ der Regulator des Systems der $\mathcal{E}_i^{(1)}$, so ergiebt sich nach dem Multiplicationssatze der Determinanten, wenn mit $C$ der absolute Werth der Determinante der ganzzahligen Exponenten $\xi_{i,k}$, also eine \gesperrt{ganze rationale positive Zahl} bezeichnet wird:

\weq{13}{T_1 = C\,T_0.}

Es kann nicht bewiesen werden und trifft für die höheren Werthe von $m$ wahrscheinlich auch nicht zu, dass $\tau_0, \tau_1, \ldots, \tau_{\nu'-1}$ ein Fundamentalsystem bilden\footnote{Es ist hier, wie bei allen Fragen, die die Einheiten betreffen, sehr schwierig, zu bestimmten numerischen Resultaten zu gelangen. Nach den Resultaten, die in dem reichhaltigen Werke von \gesperrt{Reuschle}, \glqq Tafeln complexer Primzahlen, welche aus Wurzeln der Einheiten gebildet sind\grqq{} (Berlin 1875), enthalten sind, bilden für $m=16$ die $\tau$ noch ein Fundamentalsystem. Denn hier ist die Grundzahl des Körpers $H_m$ gleich $2^{11}$, der Grad des Körpers gleich 4, und daher giebt es nach der Anmerkung zu \S.\,156 in jeder Idealclasse ein Ideal, dessen Norm kleiner als 6 ist. Nach den Tafeln von \gesperrt{Reuschle} kann aber jede natürliche Primzahl unter 100 im Körper $H_m$ in wirklich existirende Primfactoren zerlegt werden. Folglich gehen gewiss in allen Zahlen unter 6 nur Hauptideale auf, und folglich giebt es hier nur die eine Classe, die Hauptclasse. Da, wie wir gesehen haben, auch der erste Classenzahlfactor $=1$ ist, so ist der Kreistheilungskörper $\Omega_{16}$ einclassig. Es gelten in ihm dieselben Gesetze der Primzahlen, wie im Körper der rationalen Zahlen. Zweifelhaft ist dies für den Körper $\Omega_{32}$ und sicher nicht mehr zutreffend im Körper $\Omega_{64}$, in dem die Classenzahl ein Vielfaches von 17 ist.}.

Für unseren Zweck ist es aber von Wichtigkeit, dass sich diese Einheiten in Bezug auf den Nenner 2 in den Exponenten
%CONT>>
